%% GENERATED by tools/gen_latex.py -- edit content/ch20.py instead.
%% Build:  latexmk -pdf ch20.tex     (or: pdflatex ch20.tex, twice)
\documentclass[aspectratio=169,11pt,t]{beamer}
\usepackage{itbpro}

\coursetitle{Conclusions}
\coursesubtitle{Everything this course covered, arranged so you can reconstruct the decisions rather than remember the details --- and where to go next.}
\coursekicker{Chapter 20}
\coursesource{Chollet \& Watson, Deep Learning with Python 3e --- chapter 20}
\courseduration{2 hours (1 session)}
\coursepresenter{Prof. Bambang Riyanto Trilaksono $\cdot$ Rahman Indra Kesuma, S.Kom., M.Cs.}{Lead Instructor $\cdot$ Teaching Assistant}
\coursebrand{AI for Professional \textperiodcentered\ ITB \textperiodcentered\ Ch 20}

\title{Conclusions}
\author{Prof. Bambang Riyanto Trilaksono $\cdot$ Rahman Indra Kesuma, S.Kom., M.Cs.}
\date{}

\begin{document}
\covertitle

\framekicker{Learning outcomes}
\begin{frame}[shrink=25]{Session Objectives}
\leadpar{By the end of this session, participants can:}
\begin{isteps}
\item Place \textbf{AI, machine learning, deep learning, and generative AI} correctly inside one another.
\item State in one paragraph \textbf{what a deep learning model is} and how learning happens.
\item Name the \textbf{four enabling factors} behind the current revolution, and why no single one of them explains it.
\item Run the \textbf{universal workflow} from problem definition to deployment, and say what each step protects against.
\item Pick the right \textbf{architecture family} for a given pair of input and output modalities.
\item Write the correct \textbf{final layer, activation, and loss} for each of the five standard task types from memory.
\item State the \textbf{three structural limitations} honestly to a non-technical stakeholder.
\item Keep learning after this course: \textbf{Kaggle, arXiv, and the Keras ecosystem}.
\end{isteps}

\vskip 6pt
\vskip 4pt
\reslink{Book}{Chapter 20 --- full text}{https://deeplearningwithpython.io/chapters/chapter20\_conclusion}
\vskip 2pt
\end{frame}

\coursesection{01}{Key concepts in review}{If you ever need a quick refresher, this is the section to reread.}

\framekicker{Section 20.1.1}
\begin{frame}[shrink=25]{Four terms, nested inside one another}
{\fontsize{9.4}{12.6}\selectfont\color{fg2}Deep learning is not synonymous with artificial intelligence, or even with machine learning. The distinctions matter, and they are routinely blurred in public discussion.\par}\vskip 4pt
\begin{center}
\includegraphics[width=0.92\linewidth,height=0.60\textheight,keepaspectratio]{figs/ch20-nesting.tex.pdf}
\end{center}
\vskip 3pt{\color{ink3}\fontsize{7.4}{9.4}\selectfont Each term is a proper subset of the one above it.\par}
\end{frame}

\framekicker{Section 20.1.1 $\cdot$ the definitions}
\begin{frame}[shrink=25]{Said precisely}
\begin{itable}{>{\raggedright\arraybackslash}p{0.2256\textwidth}>{\raggedright\arraybackslash}p{0.7144\textwidth}}
\thead{Term} & \thead{Definition} \\ \midrule
\textbf{AI} & An ancient, broad field: \textbf{all attempts to automate human cognitive processes} --- from an Excel spreadsheet to a humanoid robot. \\
\textbf{Machine learning} & Automatically developing programs (\textbf{models}) purely from exposure to training data. Took off in the 1990s, dominant by the 2000s. \\
\textbf{Deep learning} & Models that are \textbf{long chains of geometric transformations}, structured into layers parameterized by weights. \\
\textbf{Generative AI} & Deep learning models that generate text, images, video, or sound --- very large, trained \textbf{self-supervised} to reconstruct corrupted input. \\
\end{itable}
{\fontsize{9.4}{12.6}\selectfont\color{fg2}The self-supervised objective is what lets generative models learn sophisticated \textbf{maps of their input space} --- embedding manifolds that can then be sampled from.\par}\vskip 4pt
\end{frame}

\framekicker{Section 20.1.2}
\begin{frame}[shrink=25]{Why deep learning is not on equal footing with the alternatives}
{\fontsize{9.4}{12.6}\selectfont\color{fg2}In a few years, deep learning achieved breakthroughs across tasks historically perceived as extremely difficult --- especially \textbf{machine perception}: extracting useful information from images, video, and sound.\par}\vskip 4pt
{\fontsize{9.4}{12.6}\selectfont\color{fg2}Given sufficient training data, appropriately labelled, deep learning can extract from perceptual data \textbf{almost anything a human could}. It is sometimes said to have \textit{solved perception} --- true only for a fairly narrow definition of perception.\par}\vskip 4pt
\begin{iband}[signal]
This singlehandedly brought about the \textbf{third and by far the largest AI summer}: a period of intense interest, investment, and hype. We are in the middle of it.
\end{iband}
\end{frame}

\framekicker{Section 20.1.2 $\cdot$ why this summer is different}
\begin{frame}[shrink=25]{The hype may recede; the impact will not}
{\fontsize{9.4}{12.6}\selectfont\color{fg2}In stark contrast with previous AI summers, deep learning has provided \textbf{enormous business value} to companies large and small, and become a consumer success: human-level speech recognition, chatbot assistants, photorealistic image generation, human-level machine translation.\par}\vskip 4pt
\begin{iquote}
\quotetext{Deep learning could be analogous to \textbf{the internet}: overly hyped for a few years, but in the longer term still a major revolution that will transform our economy and our lives.}\quotecite{Section 20.1.2}
\end{iquote}
\begin{iband}[signal]
One reason for optimism: \textbf{even with no further technological progress in the next decade, deploying existing algorithms to every applicable problem would be a game changer for most industries.} Short-term expectations are overoptimistic; the full deployment will likely take decades.
\end{iband}
\end{frame}
\note{This is the framing to give a professional audience about their own roadmap. The value is not in waiting for the next model --- it is in applying what already exists to problems nobody has applied it to yet.}

\framekicker{Section 20.1.3}
\begin{frame}[shrink=25]{The most surprising thing is how simple it is}
{\fontsize{9.4}{12.6}\selectfont\color{fg2}Fifteen years ago nobody expected such results on machine perception and natural language from \textbf{simple parametric models trained with gradient descent}. It turns out all you need is sufficiently large models, trained on sufficiently many examples.\par}\vskip 4pt
\begin{iquote}
\quotetext{\textit{"It's not complicated, it's just a lot of it."}}\quotecite{Richard Feynman, on the universe --- quoted in section 20.1.3}
\end{iquote}
\begin{center}
\includegraphics[width=0.92\linewidth,height=0.60\textheight,keepaspectratio]{figs/ch20-learning.tex.pdf}
\end{center}
\vskip 3pt{\color{ink3}\fontsize{7.4}{9.4}\selectfont The knowledge of a model is stored in its weights; learning is finding good values for them.\par}
\end{frame}

\framekicker{Section 20.1.3}
\begin{frame}[shrink=25]{In deep learning, everything is a vector}
{\fontsize{9.4}{12.6}\selectfont\color{fg2}Inputs and targets are first \textbf{vectorized} --- turned into an input vector space and a target vector space. Each layer operates one simple geometric transformation. Together the chain forms one complex transformation, broken into simple steps.\par}\vskip 4pt
{\fontsize{9.4}{12.6}\selectfont\color{fg2}That transformation must be \textbf{differentiable}, which is what lets gradient descent learn its parameters. Intuitively: the morphing from inputs to outputs must be \textbf{smooth and continuous} --- a significant constraint, and the source of every limitation in chapter 19.\par}\vskip 4pt
\begin{center}
\includegraphics[width=0.92\linewidth,height=0.60\textheight,keepaspectratio]{figs/ch20-uncrumple.tex.pdf}
\end{center}
\vskip 3pt{\color{ink3}\fontsize{7.4}{9.4}\selectfont Deep learning models are mathematical machines for uncrumpling complicated manifolds of high-dimensional data.\par}
\end{frame}

\framekicker{Section 20.1.3 $\cdot$ the two core ideas}
\begin{frame}[shrink=25]{Meaning as geometry}
\leadpar{The whole thing hinges on two ideas: that \textbf{meaning derives from the pairwise relationships between things} --- words in a language, pixels in an image --- and that \textbf{those relationships can be captured by a distance function}.}
{\fontsize{9.4}{12.6}\selectfont\color{fg2}Whether the brain implements meaning via geometric spaces is an entirely separate question. Vector spaces are computationally efficient, but other data structures for intelligence are easily envisioned --- \textbf{graphs} in particular.\par}\vskip 4pt
{\fontsize{9.4}{12.6}\selectfont\color{fg2}Neural networks originally emerged from the idea of using graphs to encode meaning; the surrounding field was called \textbf{connectionism}. That is the only reason for the name today.\par}\vskip 4pt
\end{frame}

\framekicker{Section 20.1.3 $\cdot$ a naming complaint worth taking seriously}
\begin{frame}[shrink=25]{They are neither neural nor networks}
\begin{iband}[amber]
\textit{Neural network} exists \textbf{purely for historical reasons}. It is an extremely misleading name --- these are neither neural nor networks, and they have \hilite{hardly anything to do with the brain}.
\end{iband}
{\fontsize{9.4}{12.6}\selectfont\color{fg2}More appropriate names, all of which emphasise that continuous geometric space manipulation is at the core:\par}\vskip 4pt
\begin{ibullets}
\item \textbf{Layered representations learning}
\item \textbf{Hierarchical representations learning}
\item \textbf{Deep differentiable models}
\item \textbf{Chained geometric transforms}
\end{ibullets}
{\fontsize{9.4}{12.6}\selectfont\color{fg2}This is not pedantry. The brain metaphor is precisely what invites the anthropomorphizing that chapter 19 warned against.\par}\vskip 4pt
\end{frame}

\framekicker{Section 20.1.4}
\begin{frame}[shrink=25]{Four enabling factors, no single breakthrough}
{\fontsize{9.4}{12.6}\selectfont\color{fg2}The revolution did not start with one invention. Like any revolution it is the product of an accumulation of enabling factors --- \textbf{slowly at first, and then suddenly}.\par}\vskip 4pt
\begin{center}
\includegraphics[width=0.92\linewidth,height=0.60\textheight,keepaspectratio]{figs/ch20-enablers.tex.pdf}
\end{center}
\vskip 3pt{\color{ink3}\fontsize{7.4}{9.4}\selectfont Remove any one of the four and the last decade does not happen.\par}
\end{frame}

\framekicker{Section 20.1.4 $\cdot$ two details worth keeping}
\begin{frame}[shrink=25]{Where the data and the hardware came from}
\begin{minipage}[t]{0.494\textwidth}
\begin{icardaccent}
\cardh{The data}
A by-product of the \textbf{consumer internet} and Moore's law applied to storage media. Today, state-of-the-art language models are trained on \textbf{a large fraction of the entire internet}.
\end{icardaccent}
\end{minipage}%
\hfill
\begin{minipage}[t]{0.494\textwidth}
\begin{icardaccent}
\cardh{The hardware}
First \textbf{gaming GPUs}, then chips designed from the ground up for deep learning. NVIDIA's CEO took note of the boom early and bet the company's future on it.
\end{icardaccent}
\end{minipage}%
{\fontsize{9.4}{12.6}\selectfont\color{fg2}And on top of both, a stack of software: \textbf{CUDA}, then \textbf{TensorFlow, JAX, and PyTorch} doing automatic differentiation, then \textbf{Keras} making the whole thing accessible to most people.\par}\vskip 4pt
\end{frame}

\framekicker{Section 20.1.4 $\cdot$ the direction of travel}
\begin{frame}[shrink=25]{A tool in every developer's toolbox}
\leadpar{In the future, deep learning will not be used only by researchers and engineers with an academic profile. \textbf{It will be a tool in the toolbox of every developer, much like web technology today.}}
{\fontsize{9.4}{12.6}\selectfont\color{fg2}Everyone needs to build intelligent apps. Just as every business today needs a website, \textbf{every product will need to intelligently make sense of user-generated data}.\par}\vskip 4pt
{\fontsize{9.4}{12.6}\selectfont\color{fg2}Bringing that about requires tools that make deep learning radically easy to use and accessible to anyone with basic coding ability. \hilite{Keras has been the first major step in that direction}, and it is why this course is taught in it.\par}\vskip 4pt
\end{frame}

\coursesection{02}{The universal workflow}{The hard part is everything before the model, and everything after it.}

\framekicker{Section 20.1.5}
\begin{frame}[shrink=25]{Keras cannot help you with the hard part}
{\fontsize{9.4}{12.6}\selectfont\color{fg2}Having an extremely powerful tool for mapping any input space to any target space is great. But the difficult part of the workflow is \textbf{often everything that comes before} designing and training the model --- and for production models, everything that comes after.\par}\vskip 4pt
\begin{iband}[signal]
Understanding the problem domain well enough to determine \textbf{what to predict, from what data, and how to measure success} is a prerequisite for any successful application. \hilite{It is not something advanced tools can help you with.}
\end{iband}
\end{frame}

\framekicker{Section 20.1.5 $\cdot$ from chapter 6}
\begin{frame}[shrink=25]{The workflow, end to end}
\begin{center}
\includegraphics[width=0.92\linewidth,height=0.60\textheight,keepaspectratio]{figs/ch20-workflow.tex.pdf}
\end{center}
\vskip 3pt{\color{ink3}\fontsize{7.4}{9.4}\selectfont Seven steps, of which only one is about the model.\par}
\end{frame}

\framekicker{Section 20.1.5}
\begin{frame}[shrink=25]{What each step protects against}
\begin{itable}{>{\raggedright\arraybackslash}p{0.3196\textwidth}>{\raggedright\arraybackslash}p{0.6204\textwidth}}
\thead{Step} & \thead{The failure it prevents} \\ \midrule
Define the problem & Building a model for a prediction nobody needs, or that the data cannot support. \\
Choose a measure of success & Optimizing a convenient number instead of a useful one --- the shortcut rule from chapter 19. \\
Prepare the validation process & \textbf{Label leakage.} With temporal prediction, validation and test data must be \textit{posterior} to the training data. \\
Beat a common-sense baseline & Spending months on a problem where machine learning does not work at all. \\
Overfit, then regularize & Regularizing a model that was never big enough to fit the data in the first place. \\
Keep a separate test set & \textbf{Validation-set overfitting} from hyperparameter tuning. This is the entire purpose of the third split. \\
\end{itable}
\end{frame}

\framekicker{Section 20.1.5 $\cdot$ the step people skip}
\begin{frame}[shrink=25]{Deployment is part of the workflow, not after it}
{\fontsize{9.4}{12.6}\selectfont\color{fg2}Deploy the final model in production --- as a web API, inside a JavaScript or C++ application, on an embedded device.\par}\vskip 4pt
\leadpar{Then \textbf{keep monitoring its performance on real-world data}, and use those findings to refine the next iteration.}
{\fontsize{9.4}{12.6}\selectfont\color{fg2}The arrow from deployment back to problem definition is the one that turns a project into a system. Chapter 19's warning about distribution shift is not a research concern: \textbf{production data drifts away from your training set from the day you ship.}\par}\vskip 4pt
\end{frame}

\coursesection{03}{Key network architectures}{An architecture encodes assumptions about the structure of the data.}

\framekicker{Section 20.1.6}
\begin{frame}[shrink=25]{An architecture is a hypothesis space}
\leadpar{A network architecture \textbf{encodes assumptions about the structure of the data} --- a hypothesis space within which the search for a good model proceeds.}
{\fontsize{9.4}{12.6}\selectfont\color{fg2}Whether an architecture works on a problem depends \textbf{entirely} on the match between the structure of the data and the assumptions of the architecture. That is the single question to ask when choosing one.\par}\vskip 4pt
\begin{iband}[signal]
These types combine easily into larger multimodal models. \textbf{Deep learning layers are LEGO bricks for information processing.}
\end{iband}
\end{frame}

\framekicker{Section 20.1.6 $\cdot$ table 20.1}
\begin{frame}[shrink=25]{Input, output, architecture}
\begin{itable}{>{\raggedright\arraybackslash}p{0.282\textwidth}>{\raggedright\arraybackslash}p{0.376\textwidth}>{\raggedright\arraybackslash}p{0.282\textwidth}}
\thead{Input} & \thead{Output} & \thead{Model} \\ \midrule
Vector data & Class probability, regression value & \textbf{Densely connected network} \\
Timeseries data & Class probability, regression value & \textbf{RNN, Transformer} \\
Images & Class probability, regression value & \textbf{ConvNet} \\
Text & Class probability, regression value & \textbf{Transformer} \\
Text, images & Text & \textbf{Transformer} \\
Text, images & Images & \textbf{VAE, diffusion model} \\
\end{itable}
{\fontsize{9.4}{12.6}\selectfont\color{fg2}Six rows covering the whole course. If you remember nothing else from this deck, \textbf{remember this table} --- it turns a vague problem statement into a starting architecture.\par}\vskip 4pt
\end{frame}

\framekicker{Section 20.1.6 $\cdot$ densely connected networks}
\begin{frame}[shrink=25]{No assumed structure at all}
{\fontsize{9.4}{12.6}\selectfont\color{fg2}A stack of \inlinecode{Dense} layers, for vector data where each sample is a vector of numerical or categorical attributes. They assume \textbf{no specific structure} in the input features --- every unit connects to every other, so the layer maps relationships between \textbf{any} two features.\par}\vskip 4pt
{\fontsize{9.4}{12.6}\selectfont\color{fg2}Contrast a 2D convolution layer, which only looks at \textbf{local} relationships. The absence of assumptions is both the strength and the weakness.\par}\vskip 4pt
{\fontsize{9.4}{12.6}\selectfont\color{fg2}They are also the \textbf{final classification or regression stage} of most networks: the ConvNets of chapter 8 and the recurrent networks of chapter 13 all end in one or two \inlinecode{Dense} layers.\par}\vskip 4pt
\end{frame}

\framekicker{Section 20.1.6 $\cdot$ the five endings}
\begin{frame}[shrink=25]{Last layer, activation, and loss --- from memory}
\begin{itable}{>{\raggedright\arraybackslash}p{0.3008\textwidth}>{\raggedright\arraybackslash}p{0.3196\textwidth}>{\raggedright\arraybackslash}p{0.3196\textwidth}}
\thead{Task} & \thead{Last layer} & \thead{Loss} \\ \midrule
\textbf{Binary classification} & \inlinecode{Dense(1, activation="sigmoid")} & \inlinecode{binary\_crossentropy} \\
\textbf{Single-label, multiclass} & \inlinecode{Dense(n, activation="softmax")} & \inlinecode{categorical\_crossentropy} (one-hot targets) \\
\textbf{Single-label, integer targets} & \inlinecode{Dense(n, activation="softmax")} & \inlinecode{sparse\_categorical\_crossentropy} \\
\textbf{Multilabel, multiclass} & \inlinecode{Dense(n, activation="sigmoid")} & \inlinecode{binary\_crossentropy} (k-hot targets) \\
\textbf{Regression} & \inlinecode{Dense(n)} --- \textbf{no activation} & \inlinecode{mean\_squared\_error} \\
\end{itable}
\begin{iband}[amber]
The multilabel row is the one people get wrong. \textbf{Sigmoid, not softmax} --- softmax forces the outputs to sum to one, which is exactly what you do not want when a sample can have several classes.
\end{iband}
\end{frame}

\framekicker{Section 20.1.6 $\cdot$ densely connected networks}
\begin{frame}[shrink=25]{The same three lines, three different endings}
{\fontsize{9.4}{12.6}\selectfont\color{fg2}The body of the network is identical in all three cases. Only the last layer and the loss change.\par}\vskip 4pt
\icode{python}{binary classification}{listings/ch20-01.py}
\icode{python}{single-label multiclass}{listings/ch20-02.py}
{\fontsize{9.4}{12.6}\selectfont\color{fg2}Targets should be \textbf{0 or 1} for the first, \textbf{one-hot} for the second --- or integers, with \inlinecode{sparse\_categorical\_crossentropy}.\par}\vskip 4pt
\end{frame}

\framekicker{Section 20.1.6 $\cdot$ Transformers}
\begin{frame}[shrink=25]{A sequence-to-sequence Transformer, from KerasHub layers}
{\fontsize{9.4}{12.6}\selectfont\color{fg2}Everything chapter 15 built by hand exists as a layer. This shape covers machine translation and question answering alike.\par}\vskip 4pt
\icode{python}{}{listings/ch20-03.py}
{\fontsize{9.4}{12.6}\selectfont\color{fg2}\inlinecode{TokenAndPositionEmbedding} is the tied token-plus-position layer from chapters 15 and 16. \textbf{If you are processing a single sequence, delete the decoder half} and keep the encoder.\par}\vskip 4pt
\end{frame}

\framekicker{Section 20.1.6 $\cdot$ ConvNets}
\begin{frame}[shrink=25]{The same transformation, applied everywhere}
{\fontsize{9.4}{12.6}\selectfont\color{fg2}Convolution layers look at \textbf{spatially local patterns} by applying the same geometric transformation to different patches of an input tensor. The result is \textbf{translation invariant}, which makes convolution layers highly data efficient and modular.\par}\vskip 4pt
{\fontsize{9.4}{12.6}\selectfont\color{fg2}The idea applies to any dimensionality: \inlinecode{Conv1D} for continuous sequences, \inlinecode{Conv2D} for images, \inlinecode{Conv3D} for volumes. \inlinecode{SeparableConv2D} is a leaner, more efficient alternative.\par}\vskip 4pt
{\fontsize{9.4}{12.6}\selectfont\color{fg2}ConvNets are stacks of convolution and pooling layers. Pooling \textbf{downsamples spatially} --- required to keep feature maps a reasonable size as feature counts grow, and to let later convolutions see a greater spatial extent.\par}\vskip 4pt
\end{frame}

\framekicker{Section 20.1.6 $\cdot$ ConvNets}
\begin{frame}[shrink=25]{A typical image classifier, and what to add when it goes deep}
{\fontsize{9.4}{12.6}\selectfont\color{fg2}Four separable convolutions, one pooling stage, a global pool, and a dense head.\par}\vskip 4pt
\icode{python}{}{listings/ch20-04.py}
{\fontsize{9.4}{12.6}\selectfont\color{fg2}The \inlinecode{GlobalAveragePooling2D} turns spatial feature maps into vectors, which the \inlinecode{Dense} layers then classify. When building a \textbf{very deep} ConvNet, add \textbf{batch normalization} and \textbf{residual connections} --- the two patterns from chapter 9 that keep gradient information flowing.\par}\vskip 4pt
\end{frame}

\framekicker{Section 20.1.6 $\cdot$ Transformers}
\begin{frame}[shrink=25]{Context awareness for a set of vectors}
{\fontsize{9.4}{12.6}\selectfont\color{fg2}A Transformer looks at a \textbf{set} of vectors and uses neural attention to transform each into a representation that is aware of the context provided by the others.\par}\vskip 4pt
{\fontsize{9.4}{12.6}\selectfont\color{fg2}When the set is an ordered sequence, \textbf{positional encoding} adds order awareness --- giving a model that handles both global context and word order, and processes long paragraphs far more effectively than RNNs or 1D ConvNets.\par}\vskip 4pt
\begin{minipage}[t]{0.494\textwidth}
\begin{icardaccent}
\cardh{TransformerEncoder}
Turns an input vector sequence into a \textbf{context-aware, order-aware} output sequence. If you are processing a single sequence, this is all you need.
\end{icardaccent}
\end{minipage}%
\hfill
\begin{minipage}[t]{0.494\textwidth}
\begin{icardaccent}
\cardh{TransformerDecoder}
Takes the encoder output \textbf{and} a target sequence, and predicts what comes next in the target.
\end{icardaccent}
\end{minipage}%
\end{frame}

\framekicker{Section 20.1.6 $\cdot$ the remaining families}
\begin{frame}[shrink=25]{RNNs, VAEs, and diffusion models}
\begin{minipage}[t]{0.3253\textwidth}
\begin{icard}
\cardh{Recurrent networks}
Process sequences one step at a time, carrying state forward. Still the natural fit for \textbf{timeseries} --- chapter 13 showed a small GRU beating a much larger model.
\end{icard}
\end{minipage}%
\hfill
\begin{minipage}[t]{0.3253\textwidth}
\begin{icard}
\cardh{VAEs}
Encode to a \textbf{distribution} rather than a point, giving a continuous, structured latent space. Good where control and interpretability matter more than fidelity.
\end{icard}
\end{minipage}%
\hfill
\begin{minipage}[t]{0.3253\textwidth}
\begin{icardaccent}
\cardh{Diffusion models}
A denoising autoencoder \textbf{in a loop}, turning pure noise into an image. Behind nearly every commercial image generator today.
\end{icardaccent}
\end{minipage}%
{\fontsize{9.4}{12.6}\selectfont\color{fg2}All of these combine. A multimodal model splices an image encoder's output into a text sequence; a text-to-image model feeds a Transformer's embeddings to a diffusion denoiser. \textbf{The bricks fit together.}\par}\vskip 4pt
\end{frame}

\coursesection{04}{Limitations, and what might lie ahead}{State these plainly. They are the most useful thing you can tell a stakeholder.}

\framekicker{Section 20.2}
\begin{frame}[shrink=25]{Strength and weakness are the same property}
{\fontsize{9.4}{12.6}\selectfont\color{fg2}Layers plug together to map essentially anything to anything --- given appropriate training data, and given that the mapping is achievable by a \textbf{continuous geometric transformation of reasonable complexity}.\par}\vskip 4pt
\begin{center}
\includegraphics[width=0.92\linewidth,height=0.60\textheight,keepaspectratio]{figs/ch20-limits.tex.pdf}
\end{center}
\vskip 3pt{\color{ink3}\fontsize{7.4}{9.4}\selectfont Here is the catch: the mapping is often not learnable in a way that generalizes.\par}
\end{frame}

\framekicker{Section 20.2 $\cdot$ the three limits, for a stakeholder}
\begin{frame}[shrink=25]{How to say this without hedging}
\begin{isteps}
\item \textbf{They struggle to adapt to novelty.} Parameters are fixed after training, so the model can only retrieve or replicate patterns similar to its training data. Outside that distribution, performance degrades drastically --- \textit{no matter how simple the underlying task}.
\item \textbf{They are sensitive to phrasing and other distractors.} Minor rewording in a prompt, or imperceptible perturbation in an image, changes the answer --- which indicates the absence of robust, human-like understanding.
\item \textbf{They often cannot learn generalizable algorithms.} The continuous, geometric nature of these models makes them ill-suited to exact, discrete, step-by-step procedures. They approximate such processes by interpolation.
\end{isteps}
\begin{iband}[rose]
And always resist the temptation to anthropomorphize. Performance is built on \textbf{pointwise statistical patterns}, not human-like experiential grounding --- which is why it is brittle at any deviation from the training data.
\end{iband}
\end{frame}
\note{This slide is the one to rehearse. Being able to state these three limits calmly, without either overselling or dismissing the technology, is the most valuable professional skill this course teaches.}

\framekicker{Section 20.2 $\cdot$ where scaling got us}
\begin{frame}[shrink=25]{Five years of exponential scaling, and an open problem}
{\fontsize{9.4}{12.6}\selectfont\color{fg2}The narrative that scaling model size and data would produce general intelligence has proven \textbf{insufficient}. Scaling improves benchmarks that amount to memorization tests; it does not address limits that stem from \textbf{fitting static, interpolative curves to data}.\par}\vskip 4pt
{\fontsize{9.4}{12.6}\selectfont\color{fg2}By 2024 this spurred the transition to \textbf{test-time adaptation} --- search or fine-tuning during inference. It yielded major breakthroughs, including o3 surpassing the human baseline on ARC-AGI-1.\par}\vskip 4pt
\begin{iband}[amber]
But \textbf{at extreme computational cost}. Efficient, human-like adaptation remains a completely open problem, and the slightly harder \textbf{ARC-AGI-2 remains completely unsolved}. \hilite{We still need conceptual advances beyond scaling or brute-force search.}
\end{iband}
\end{frame}

\framekicker{Section 20.3}
\begin{frame}[shrink=25]{Three developments worth watching}
\begin{center}
\includegraphics[width=0.92\linewidth,height=0.60\textheight,keepaspectratio]{figs/ch20-ahead.tex.pdf}
\end{center}
\vskip 3pt{\color{ink3}\fontsize{7.4}{9.4}\selectfont Deep learning excels at value-centric abstraction and lacks program-centric abstraction. Human intelligence integrates both.\par}
{\fontsize{9.4}{12.6}\selectfont\color{fg2}The last of the three is the one with the clearest analogy to software engineering: \textbf{libraries of reusable, modular components} acquired from experience, fetched and recombined into a model adapted to the situation at hand.\par}\vskip 4pt
\end{frame}

\coursesection{05}{Staying up to date}{What you learned here will not stay relevant forever.}

\framekicker{Section 20.4}
\begin{frame}[shrink=25]{The field moves at a frenetic pace}
{\fontsize{9.4}{12.6}\selectfont\color{fg2}Modern deep learning is only a few years old, despite a long, slow prehistory stretching back decades. With exponentially increasing financial resources and research headcount since 2013, the field is now moving very fast.\par}\vskip 4pt
\leadpar{\textbf{What you have learned in this course will not stay relevant forever, and it is not all you will need for the rest of your career.}}
\begin{center}
\includegraphics[width=0.92\linewidth,height=0.60\textheight,keepaspectratio]{figs/ch20-keepgoing.tex.pdf}
\end{center}
\vskip 3pt{\color{ink3}\fontsize{7.4}{9.4}\selectfont Three free resources, and what each one is actually good for.\par}
\end{frame}

\framekicker{Section 20.4.1}
\begin{frame}[shrink=25]{Kaggle: the only real way to learn is practice}
{\fontsize{9.4}{12.6}\selectfont\color{fg2}An array of constantly renewed data science competitions, many involving deep learning, prepared by companies wanting novel solutions to their hardest problems. Fairly large monetary prizes for top entrants.\par}\vskip 4pt
{\fontsize{9.4}{12.6}\selectfont\color{fg2}Participating in a few --- perhaps as part of a team --- makes you familiar with the practical side of the advanced practices in chapter 18:\par}\vskip 4pt
\begin{ibullets}
\item \textbf{Hyperparameter tuning} at a scale where it actually matters.
\item \textbf{Avoiding validation-set overfitting} when the leaderboard is watching.
\item \textbf{Model ensembling}, and finding out first-hand that diversity beats individual quality.
\end{ibullets}
\end{frame}

\framekicker{Section 20.4.2}
\begin{frame}[shrink=25]{arXiv: the field publishes in the open}
{\fontsize{9.4}{12.6}\selectfont\color{fg2}Deep learning research takes place \textbf{completely in the open}. Papers are freely accessible as soon as they are finalized, and much of the related software is open source.\par}\vskip 4pt
{\fontsize{9.4}{12.6}\selectfont\color{fg2}Researchers upload to arXiv shortly after completion, to plant a flag without waiting months for conference acceptance. \textbf{This is what lets the field move so fast} --- every new finding is immediately available to build on.\par}\vskip 4pt
\begin{iband}[amber]
The downside: the sheer volume makes it impossible to skim, and papers are \textbf{not peer-reviewed}, so identifying what is both important and high quality is hard and getting harder. \hilite{Use Google Scholar to track publications by authors you trust.}
\end{iband}
\end{frame}

\framekicker{Section 20.4.3}
\begin{frame}[shrink=25]{The Keras ecosystem}
{\fontsize{9.4}{12.6}\selectfont\color{fg2}With over \textbf{2.5 million users} as of early 2025 and still growing, Keras has a large ecosystem of tutorials, guides, and open source projects.\par}\vskip 4pt
\vskip 4pt
\reslink{DOCS}{keras.io --- the main reference}{https://keras.io}
\reslink{GUIDES}{keras.io/guides --- extensive developer guides}{https://keras.io/guides}
\reslink{EXAMPLES}{keras.io/examples --- dozens of high-quality code examples}{https://keras.io/examples}
\reslink{SOURCE}{github.com/keras-team/keras}{https://github.com/keras-team/keras}
\reslink{HUB}{github.com/keras-team/keras-hub}{https://github.com/keras-team/keras-hub}
\vskip 2pt
{\fontsize{9.4}{12.6}\selectfont\color{fg2}The examples in particular are worth reading before you write anything from scratch --- most of what you need has a \textbf{maintained, tested reference implementation} there already.\par}\vskip 4pt
\end{frame}

\framekicker{The course, in one view}
\begin{frame}[shrink=25]{Twenty chapters, six ideas}
\begin{itable}{>{\raggedright\arraybackslash}p{0.1504\textwidth}>{\raggedright\arraybackslash}p{0.7896\textwidth}}
\thead{Chapters} & \thead{The idea} \\ \midrule
1-3 & \textbf{Learning is gradient descent on a differentiable chain of geometric transformations.} Everything else is engineering. \\
4-7 & \textbf{Generalization is the whole problem}, and the workflow exists to protect it --- splits, baselines, overfit-then-regularize. \\
8-12 & \textbf{Structure in the data buys you assumptions in the architecture.} Convolution is translation invariance made into a layer. \\
\end{itable}
\end{frame}

\framekicker{The course, in one view}
\begin{frame}[shrink=25]{Twenty chapters, six ideas (2 of 2)}
\begin{itable}{>{\raggedright\arraybackslash}p{0.1504\textwidth}>{\raggedright\arraybackslash}p{0.7896\textwidth}}
\thead{Chapters} & \thead{The idea} \\ \midrule
13-14 & \textbf{Order is a design decision.} Set or sequence is the question every architecture answers differently. \\
15-17 & \textbf{Attention turns correlation into proximity}, and interpolative spaces generate both generalization and hallucination. \\
18-20 & \textbf{Scale, precision, and honesty.} What each lever costs, and what none of them fixes. \\
\end{itable}
{\fontsize{9.4}{12.6}\selectfont\color{fg2}The through-line from chapter 5 to chapter 19 is a single claim, restated at increasing scale: \textbf{these models interpolate}. Everything they do well and everything they do badly follows from that.\par}\vskip 4pt
\end{frame}

\framekicker{Common failure modes}
\begin{frame}[shrink=25]{Four mistakes that survive the whole course}
\begin{minipage}[t]{0.494\textwidth}
\begin{icardbad}
\cardh{Label leakage into validation}
Random splits on temporal data; preprocessing fitted on the full dataset; \inlinecode{adapt()} called before the split. Every metric afterwards is fiction.
\end{icardbad}
\end{minipage}%
\hfill
\begin{minipage}[t]{0.494\textwidth}
\begin{icardbad}
\cardh{Regularizing before overfitting}
You cannot tell whether a model is too small or too regularized until you have seen it \textbf{overfit once}. Get there first.
\end{icardbad}
\end{minipage}%


\vskip 4pt

\begin{minipage}[t]{0.494\textwidth}
\begin{icardwarn}
\cardh{A convenient metric, not a useful one}
The shortcut rule applies to your project too. Netflix never shipped the model that won its own competition.
\end{icardwarn}
\end{minipage}%
\hfill
\begin{minipage}[t]{0.494\textwidth}
\begin{icardwarn}
\cardh{Anthropomorphizing the model}
Believing the model \textit{understands} leads directly to deploying it where its distribution does not reach --- and to being surprised.
\end{icardwarn}
\end{minipage}%
\end{frame}

\framekicker{Applying this after the course}
\begin{frame}[shrink=25]{Where to start, in your own organization}
\begin{isteps}
\item \textbf{Find a problem where the data already exists.} The workflow's first step is the one most projects skip, and data collection is what kills most of them.
\item \textbf{Write down the measure of success before the model.} If you cannot state it in one sentence, the project is not ready.
\item \textbf{Build the common-sense baseline first.} It costs an afternoon and occasionally ends the project --- which is a good outcome, found early.
\item \textbf{Prefer a pretrained backbone.} Chapters 8, 15, and 16 all made the same point: fine-tuning beats training from scratch, decisively.
\item \textbf{Plan for monitoring on day one.} Distribution shift is not a research concern; it starts the day you deploy.
\end{isteps}
\end{frame}
\note{Close the technical part of the course here. This is the slide participants will photograph.}

\framekicker{Section 20.5 $\cdot$ final words}
\begin{frame}[shrink=25]{Go on learning, questioning, and researching}
\begin{iquote}
\quotetext{Learning is a lifelong journey, especially in the field of AI, where we have far more unknowns on our hands than certitudes. So please go on learning, questioning, and researching. \textbf{Never stop.} Because even given the progress made so far, most of the fundamental questions in AI remain unanswered. \textbf{Many haven't even been properly asked yet.}}\quotecite{Chollet \& Watson, Deep Learning with Python, closing words}
\end{iquote}
\vskip 4pt
\reslink{BOOK}{deeplearningwithpython.io}{https://deeplearningwithpython.io/}
\reslink{MODULE}{Agentic AI --- cases, tech stack, and demos}{https://hendrikarisma.my.id/rs/ai-products-course/hendri-agentic/index.html}
\vskip 2pt
\end{frame}

\end{document}
